\scanpage{371}\mspage{287}
de $M_\pi$. On sait \deleted{que \unsure{al.} alors que}\inserted{donc que, $S(\pi)$ étant karoubienne,}
$S(\pi)\unsure{-|\longrightarrow}\bigl(\underbrace{\mathrm{Ind}_\pi(S(\pi))}\bigr)_\pi$ est une équivalence
de catégories. Comme $\mathrm{Ind}_\pi(S(\pi))\approx\underline S$,
\deleted{\unsure{al|et}}\unsure{et} que $S(\pi)$ est str. pleine dans
$M_\pi$, donc aussi dans $M$ et a fortiori
dans $\underline S$, que l'on a
\[
 \begin{aligned}
 \underline S(\pi)=\underline S_\pi,\quad&
 \inserted{\text{sous-catégorie de }\underline S\text{ des}}\ \text{objets }\pi\text{-accessibles}\\
 &\text{de }S,
 \end{aligned}
 \tag*{\unsure{(4.13 1)|(4.13.1)}}
\]
Je dis qu'on aura aussi
\[
 \underline S(\pi)=\underline S\cap M_\pi.\tag*{(4.13.1)}
\]
Comme $\subset$ est clair, il reste à voir
que $\underline S\cap M_\pi\subset\underline S(\pi)$, or il est
immédiat que $\underline S\cap M_\pi\subset S_\pi$, et
on conclut par \unsure{(4,13.1)|(4.13.1)}. On \unsure{trouve|trouvera}
ainsi\unsure{:}

\uline{Proposition} 4.13.4. Soit\deleted{\unsure{nt}} $M$ une
catégorie accessible, \deleted{$\pi_0$ le plus petit
cardinal infini utile pour $M$,} $\pi$ un
cardinal \deleted{$\geq\pi_0$}\inserted{infini}. Soit $\Sigma_\pi$ l'ens des
% archive sheet number: 370
